% Working outline: old prose cleared on 2026-11-09.
\section{RQ3 --- How much granularity is enough?}
\label{sec:rq3}

\subsection{Why this needs a new criterion}

\subsection{Building the coarse--fine hierarchy}
\label{sec:rq3-hierarchy}

For this analysis, LDA supplies a fixed coarse reference and BERTopic supplies the fine
topics. The models are fitted on the same sample for each language. The script
\texttt{build\_topic\_hierarchy.py} selects a candidate parent for each fine topic: the coarse
topic containing the largest share of its documents. It checks this share and the lead over
the next-best parent. It also checks whether centroid similarity favours the same parent and
whether the child is among that parent's largest document contributors.

A link is rejected when overlap is weak, the candidate parents are nearly tied, or the
centroid and reverse checks jointly oppose the candidate. A confidence score records
agreement between these signals. The children counted below are those with accepted links;
the rules are held fixed across the sweep. The saved mappings are in
\texttt{output/hierarchy\_probe/}. Rejecting a link leaves a fine topic without an accepted
coarse parent; it does not remove documents or change their production labels.

\subsection{The floor: a split needs at least two children}

\subsection{The sweep}

\subsection{What the sweep shows}

\subsection{Why the floor alone cannot select a setting}
\label{sec:rq3-floorlimit}

\subsection{The ceiling: do the children differ from one another?}
\label{sec:rq3-ceiling}

\subsection{Both bounds at once: counting real splits}
\label{sec:rq3-joint}

\subsection{Which of the two measures actually locates a setting}
\label{sec:rq3-null}

\subsection{What this does to the German answer}
\label{sec:rq3-german}

\subsection{Answer to RQ3}

